%% kapa: exam=1 round=28 year=2017 period=1 form=A subject=real_estate subject_ko=부동산학원론 seq=32 number=112
%% source: Q-Net 감정평가사 기출문제 (cst00309) artlSeq=5199603 file=extracted/1008/5199603/2017년 제28회 감정평가사 1차시험 문제지/1교시 A형.pdf
%% conversion: kordoc 4.24.51 -> tools/gen1.py (auto); review=proofread-sonnet; flags=none
\begin{question}{112}
부동산 보유과세와 관련된 내용으로 옳지 않은 것은?
\begin{choices}
  \item 종합부동산세는 인별 과세이고 누진세율을 채택하고 있다.
  \item 토지에 대한 종합부동산세는 종합합산과세대상인 경우에는 국내에 소재하는 해당 과세대상토지의 공시가격을 합한 금액이 3억원을 초과하는 자는 종합부동산세를 납부할 의무가 있다.
  \item 종합부동산세는 조세부담의 형평성을 제고하고 가격안정을 도모하기 위해 도입되었다.
  \item 종합부동산세는 주택에 대한 종합부동산세와 토지에 대한 종합부동산세의 세액을 합한 금액을 그 세액으로 한다.
  \item 종합부동산세의 과세기준일은 재산세의 과세기준일로 한다.
\end{choices}
\end{question}
\answer{2}
